\section{AI Agent Identity：Delegation 而不是共享秘密}

课堂在 2025 年已经指出 agent framework 的安全缺口：脚本和 service account 早已代表人或系统执行任务，LLM agent 只是让行为更动态、工具更多、自然语言意图更难精确约束。核心问题不是“agent 有没有用户名”，而是谁授权它、它代表谁、能访问哪个 audience、可执行哪些 action、何时过期以及如何撤销。

\subsection{Principal、delegation 与 OAuth}

\term{service account} 是供软件或 workload 使用的非人类账户；\term{non-human identity}（NHI）泛指 service、workload、device 或 agent principal。\term{OAuth 2.0} 是 delegated authorization framework，让 client 获得 scoped access token，而不必持有用户密码；\term{OpenID Connect}（OIDC）在 OAuth 之上提供身份层，用 ID token 表达认证结果。

\lecturefigure{12-agent-delegation.png}{Agent 需要独立 principal 和受约束 delegation，而不是继承用户或机器上的全部秘密。}{本地字幕 24:10--27:50；OAuth token exchange；概念重绘。}

\begin{knowledgebox}{读图：区分“agent 是谁”和“它代表谁”}
User/owner 提供 intent 与 consent；agent principal 证明 workload identity；authorization server 依据 delegation、scope、audience 和 expiry 签发 token；tool/resource server 再做 policy check 并写 audit。\term{token exchange} 允许系统把已有 token 交换为面向特定 audience、权限更窄的新 token，同时保留 actor 与 subject 关系，避免把用户的 broad token 直接交给 agent。
\end{knowledgebox}

\teachervoice{McKinnon 提醒学生，不要被“agent 是全新事物”迷惑：cron job、script 和 service account 早就在替系统行动。真正新增的是能力和自主性扩大，而很多 framework 仍把 security 放在最后，甚至让开发机积累可访问所有系统的 token。}

\subsection{Token 是有生命周期的 capability}

\term{access token} 是向 resource server 表达授权的凭证；bearer token 被谁拿到谁就能使用，因此必须限制 scope、audience 和 lifetime，并尽量采用 sender-constrained/bound token。Agent task 完成、owner 离职、device 风险、模型被禁用或 tool policy 变化时，都应能 revoke 现有 session/token。

\lecturefigure{13-token-lifecycle.png}{Token 是临时访问能力，需要 issue、bind、use、revoke/retire 的完整生命周期。}{OAuth/OIDC/WebAuthn 概念；概念重绘。}

\begin{knowledgebox}{读图：Least privilege 要具体到任务}
Issue 阶段确认 principal、audience、scope 和 expiry；bind 阶段关联 key/device/session，减少 token theft 后的可用性；use 阶段 resource server 检查 action 与 trace ID；revoke/retire 阶段响应 risk、job completion 或 owner removal。Agent 若只需读取一个 issue 并创建草稿，就不应获得 repo admin、billing 或长期 refresh token。
\end{knowledgebox}

\begin{importantbox}{Delegation record 的最小字段}
每次 agent delegation 至少记录：human/organizational owner、agent principal、subject、requested task、resource audience、approved scopes、policy/version、issue/expiry、token/session ID、tool actions、result、revocation reason。这样 incident response 才能回答“哪个人授权了哪个 agent 对哪个系统做了什么”。
\end{importantbox}

\begin{warningbox}{Agent memory 不能成为 credential warehouse}
Prompt、memory store、trace 和 local workspace 常被设计为可检索上下文，若把 API key、refresh token 或 session cookie 混入其中，模型输出、插件和调试日志都可能扩大泄露面。Secrets 应放在专门 vault，通过 short-lived credential broker 按 task 注入，并从模型上下文和普通日志中排除。
\end{warningbox}

\subsection{本章小结}

Agent identity 的关键是 actor/subject 分离、delegation、least privilege、short-lived token 和 audit。共享用户密码或长期 token 会把一次 agent/机器 compromise 变成跨系统持久访问。

